\chapter*{本册结构与学习路线}
\addcontentsline{toc}{chapter}{本册结构与学习路线}
\label{ch:plan}

\section*{这本册子解决什么问题}
\addcontentsline{toc}{section}{这本册子解决什么问题}

《\term{编曲（arrangement）}{arrangement}与配器手册》回答一个具体问题：\textbf{手里只有旋律、和弦或一段鼓点，怎样把它组织成一首完整、有层次、能在夜半里落地的作品？}它讲的是“写什么、给谁写、写在哪一层”，不是“把已经录好的声音修得多干净”。后者属于混音与母带分册。

本册默认读者已经能在夜半里新建工程、建轨、放片段、打开钢琴卷帘。若还没有，请先读《用户指南》的“工程”“轨道与片段”“钢琴卷帘”三章。

\section*{十一章地图}
\addcontentsline{toc}{section}{十一章地图}

\begin{longtable}{p{0.06\textwidth}p{0.30\textwidth}p{0.30\textwidth}p{0.25\textwidth}}
\toprule
章 & 主题 & 学完能做什么 & 夜半落点 \\
\midrule
\endfirsthead
\toprule
章 & 主题 & 学完能做什么 & 夜半落点 \\
\midrule
\endhead
1 & 编曲总览 & 用四层角色卡判断一首歌缺什么 & 六轨演示工程、独奏/静音、钢琴卷帘 \\
2 & 歌曲结构 & 画出主歌—副歌—桥段的能量曲线 & 段落标记、编曲区、总览面板 \\
3 & 节奏组 & 让鼓、贝斯、节奏吉他/键盘互相让位 & 音频轨、MIDI 轨、力度与量化 \\
4 & 和声层 & 写铺底、琶音与和弦乐器而不糊 & 和弦轨、钢琴卷帘、设备机架 \\
5 & 旋律与对位 & 写主旋律、副旋律与呼应句 & MIDI 片段复制、移调、力度泳道 \\
6 & 弦乐配器 & 把弦乐组分成合理声部 & 弦乐轨、SFZ 音色、自动化 \\
7 & 木管与铜管 & 选音区、配组合、避免压过主唱 & 多轨分层、声像与推子 \\
8 & 电子与流行 & 用合成器层次、采样与循环做质感 & 合成器参数、素材池、片段循环 \\
9 & 人声编排 & 安排主唱、和声与叠唱 & 多条人声轨、片段对齐、导出分轨 \\
10 & 实作走查 & 从民谣、流行两首示范走完整流程 & 工程保存、导出 WAV、导出分轨 \\
11 & 曲风与配器 & 按八种真实曲风的骨架独立成段，并用四把尺判断编制 & 建轨/段落标记、钢琴卷帘、自动化、Solo/Mute \\
\bottomrule
\end{longtable}

\section*{三条学习路线}
\addcontentsline{toc}{section}{三条学习路线}

\begin{itemize}
\item \textbf{零基础路线}：第 1—5 章按顺序读，每章至少做一次八小节练习；第 10 章跟着走查做一首完整的短歌；第 11 章选一种不熟悉的曲风，按“夜半骨架”再做八小节。
\item \textbf{制作人路线}：第 2、3、4、8、10、11 章优先；第 6、7 章按需查阅，写弦乐或管乐时再回来。
\item \textbf{配器路线}：第 1、2、4、6、7、9、11 章精读；第 3、5 章作为节奏与旋律的配套。
\end{itemize}

\section*{写作与验收约定}
\addcontentsline{toc}{section}{写作与验收约定}

\begin{itemize}
\item 术语第一次出现时用“中文（English）”标出，并在书末索引中登记；本册新增术语见下一页的术语表。
\item 界面截图统一用占位命令 \verb|\screenshot{图名}{说明}|。图名用小写英文连字符，管理侧按回执清单统一拍摄后回填。
\item 概念图、流程图、谱例图用 TikZ 随源交付。图中的文字只解释关系，不冒充真实界面。
\item 每章末尾的“帮助主题映射”表把本章小节映射到应用内帮助锚点；锚点用小写英文连字符，供管理侧嵌入程序帮助。
\item 涉及夜半功能时，只写活尖端代码与现有手册里真实存在的入口；未实现或计划中的功能不写，或明确标“尚未提供”。
\end{itemize}

\section*{本册术语表}
\addcontentsline{toc}{section}{本册术语表}

下表列出本册专门术语及在本册中的含义。词条已合并进共用术语表（\texttt{common/\allowbreak glossary-zh.tex}）；合并时跳过的重复与冲突见 \texttt{GLOSSARY-MERGE.md}。

\begin{longtable}{p{0.22\textwidth}p{0.28\textwidth}p{0.42\textwidth}}
\toprule
中文 & English & 在本册中的含义 \\
\midrule
\endfirsthead
\toprule
中文 & English & 在本册中的含义 \\
\midrule
\endhead
编曲层次 & arrangement layer & 节奏、和声、旋律、音色四类功能层，不是四条固定轨道。 \\
角色卡 & role card & 写任何声部前先写清“功能、音区、节奏性格”三格。 \\
能量曲线 & energy curve & 以小节为横轴，标出各段响度、密度、乐器数的走势。 \\
织体 & texture & 各声部在音区、节奏密度与音色上的组合方式。 \\
声部避让 & voice avoidance & 让两个声部不在同一音区、同一拍点争同一功能。 \\
节奏组 & rhythm section & 鼓、贝斯、节奏吉他/键盘等负责律动与根盘的声部。 \\
铺底 & pad bed & 用长音或慢速和弦垫住和声背景的写法。 \\
琶音 & arpeggio & 把和弦音按顺序拆开演奏的织体。 \\
对位 & counterpoint & 两条以上旋律线同时进行、彼此独立又协调的写法。 \\
呼应句 & call-and-response & 主句之后由另一声部回答的短句结构。 \\
弦乐组 & string section & 第一/第二小提琴、中提琴、大提琴、低音提琴的分组。 \\
木管组 & woodwind section & 长笛、双簧管、单簧管、巴松等木管乐器组合。 \\
铜管组 & brass section & 小号、圆号、长号、大号等铜管乐器组合。 \\
叠唱 & vocal stacking & 同一歌词录多遍叠成厚人声的编排手法。 \\
走查 & walkthrough & 按一首完整作品的段落顺序，逐段检查编曲与夜半工程。 \\
主歌 & verse & 讲故事、给信息的段落，能量通常较低。 \\
副歌 & chorus & 全曲最容易被记住、和声与配器最完整的段落。 \\
预副歌 & pre-chorus & 主歌与副歌之间的抬升段，任务是制造期待。 \\
桥段 & bridge & 中段对比段，常换和弦/律动/音色，不等于第二个主歌。 \\
尾奏 & outro & 收束或淡出的结束段。 \\
段落标记 & section marker & 编曲区标尺上标出段落起止的记号。 \\
配器预算 & instrumentation budget & 按段落填写“每层加/减”的表，控制能量与差异。 \\
过渡 & transition & 两段之间的连接手法：硬切、加花、和声过渡、留白。 \\
硬切 & hard cut & 在小节线或强拍直接换段，不做过门。 \\
加花 & fill & 一小节鼓/贝斯/吉他的填充，把耳朵送进下一段。 \\
律动 & groove & 重音、留空与推进共同形成的节奏性格。 \\
锚点 & anchor & 把拍子与根音钉住的声部/落点，通常是底鼓与贝斯根音。 \\
反拍 & backbeat & 军鼓在 2、4 拍等反强拍位置的重音。 \\
半拍 & half-time & 把整体时间感拉宽一倍的鼓组写法。 \\
四拍四 & four-on-the-floor & 底鼓每拍都落的舞曲型骨架。 \\
根音 & root & 和弦在低音上的基准音；贝斯首先要把根音说清。 \\
同音重复率 & same-note repetition rate & 贝斯/律动声部中重复同一音高的比例。 \\
大跳 & leap & 相邻音之间较大的音程跳进，本册以 5 半音及以上计。 \\
压缩伴奏 & comping & 节奏吉他/键盘在关键拍给和弦、其余拍留空的伴奏写法。 \\
量化网格 & quantize grid & 音符对齐的时间格；骨架对齐它，装饰音再做人性化。 \\
和声层 & harmony layer & 承担和弦走向、织体与和声支撑的功能层。 \\
和弦节奏 & harmonic rhythm & 和弦变化的速率，如每小节/每两拍一次。 \\
声部连接 & voice leading & 相邻和弦之间各声部怎样移动，优先平滑。 \\
密集排列 & close voicing & 和弦音尽量靠近的排列，适合中高频。 \\
开放排列 & open voicing & 和弦音拉开的排列，适合铺底与弦乐。 \\
转位 & inversion & 把根音从最低声部移开，用三/五/七音做低音。 \\
顶音 & top voice & 和弦排列中最高的一声，常作为旋律线控制。 \\
和弦轨 & chord lane & 编曲区下方显示/对齐和弦变化的轨道。 \\
张力 & tension & 和弦中未解决的不稳定感，常由七/九音或延迟解决产生。 \\
持续音 & pedal & 在和弦变化时保持同一个音，制造连接或悬念。 \\
动机 & motif & 旋律里最短的可识别单位，通常 2—5 个音。 \\
乐句 & phrase & 有起有落的旋律单位，常见 4 或 8 小节。 \\
前句 & antecedent & 提出问题、建立期待的前半句。 \\
后句 & consequent & 回答前句、解决期待的后半句。 \\
轮廓 & contour & 旋律在音高上的走势：上行、下行、拱形等。 \\
目标音 & target note & 乐句落到的和弦音，决定稳定/悬/解决。 \\
副旋律 & counter-melody & 与主旋律并行、音区/节奏错开的独立旋律线。 \\
变奏 & variation & 保留动机特征、改变移位/节奏/方向的写法。 \\
呼吸 & breath & 乐句之间留出的停顿，让句子边界可听。 \\
重音 & stress & 歌词或旋律中被强调的音节/拍点。 \\
弦乐配器 & string orchestration & 把弦乐组按角色与音区分成五个声部的写法。 \\
分声部 & divisi & 一个弦乐声部内部再分成两半，以丰富和声。 \\
弓法 & bowing & 连奏、拨奏、短音等运弓方式，决定乐句性格。 \\
连奏 & legato & 音符之间不断开的演奏/写法，适合铺底与长句。 \\
拨奏 & pizzicato & 短促、颗粒清楚的弦乐写法，适合轻快段落。 \\
全奏 & tutti & 弦乐组五个声部全部发声的密度。 \\
室内乐式 & chamber-like & 只留三声部或更少的稀疏弦乐写法。 \\
低音锚 & bass anchor & 由大提琴/低音提琴守住的根音与强拍落点。 \\
声部交叉 & voice crossing & 低声部跑到高声部上方，导致角色难辨。 \\
排列 & voicing & 和弦音在各声部之间的分配与间距。 \\
木管配器 & woodwind orchestration & 用长笛、双簧管、单簧管、巴松分线写作的配器。 \\
铜管配器 & brass orchestration & 用小号、圆号、长号、大号分层与重音的配器。 \\
音色色彩 & timbral colour & 乐器音色带来的明暗、冷暖与情绪倾向。 \\
穿透力 & projection & 同等音量下乐器“钻”出混音的能力，不等于推子大小。 \\
铜管重音 & brass stab & 段落转折或副歌第一拍的短铜管和弦。 \\
同族组合 & family grouping & 木管配木管、铜管配铜管的组合方式。 \\
异族组合 & mixed grouping & 木管与铜管混用、以对比为主的组合方式。 \\
声部加倍 & doubling & 让另一件乐器重复主旋律/声部，需换音区或音程。 \\
木管二重奏 & woodwind duet & 两件木管各吹一条独立线的写法。 \\
铜管和声 & brass choir & 铜管四件按高—中—中低—低排列的和声块。 \\
电子音色层 & electronic timbre layer & 把电子音色拆成低频、主体、空气、节奏四类功能层。 \\
低频层 & sub layer & 承担最低频段与重量的合成器层。 \\
主体层 & body layer & 决定音色身份的中频合成器层。 \\
空气层 & air layer & 负责高频细节、光泽与句尾装饰的层。 \\
音色分层 & timbre layering & 让各层有独立包络与时间形状的写法。 \\
频段分工 & frequency-band division & 低频、主体、空气、节奏各占频段，不互相遮住。 \\
包络 & envelope & Attack/Decay/Release 等决定音头与尾巴的参数。 \\
滤波器 & filter & 用 Cutoff/Resonance 控制音色亮暗与峰值的设备。 \\
素材池 & asset pool & 导入音频素材的存放处，再摆到音频轨使用。 \\
采样 & sample & 导入并摆到音频轨上的现成录音材料。 \\
人声编排 & vocal arrangement & 把主唱、双唱、和声、叠唱、装饰句分角色安排的写法。 \\
主唱 & lead vocal & 唱主歌词、承担叙事与情绪的人声层。 \\
双唱 & double & 同一句歌词再唱一遍叠在主唱上的写法。 \\
人声和声 & vocal harmony & 在主唱下方三度/六度唱第二条旋律线的人声。 \\
装饰句 & ad-lib & 主唱句尾的短回应、喊唱或气声。 \\
人声层 & vocal layer & 人声编排中的一条功能层。 \\
叠唱层数 & stack layer count & 同一句歌词叠唱的层数，通常三到五层。 \\
歌词清晰度 & lyric intelligibility & 听众能否听清每个字，人声编排的首要判据。 \\
主唱中心 & lead-centered & 以主唱为轴、其他层向两侧铺开的摆位原则。 \\
分轨导出 & stem export & 把每条轨道单独导出 WAV 的交付方式。 \\
短歌 & short song & 十六小节左右的教学示范曲，不是商业长度。 \\
示范工程 & demo project & 按本册方法搭出的教学工程，不是仓库预装工程。 \\
走查清单 & walkthrough checklist & 逐段检查进出、能量、低频、主唱、色彩、交付的清单。 \\
轨道表 & track chart & 写清轨、类型与角色分配的表。 \\
段落表 & section chart & 写清小节、段落与编曲动作的表。 \\
导出回执 & export receipt & 导出 WAV/分轨后报告帧数、真峰值、条数与跳过项的回执。 \\
原子写入 & atomic write & 保存一半断电也不留下半个工程的写入方式。 \\
版本时光机 & version time machine & 把工程提交链画成图谱、可回到任意版本的面板。 \\
母线限制器 & bus limiter & 导出 WAV 时在总线上做峰值限制、防削波的开关。 \\
诊断报告包 & diagnostic report bundle & 排障用资料包，不是音频作品。 \\
曲风 & genre & 一套编制、律动、和声与结构的风格选择，不是音色标签。 \\
编制 & instrumentation & 一首曲子里用哪些乐器/声部，以及它们怎样分工。 \\
律动型 & groove pattern & 鼓、贝斯与伴奏重复的时间骨架。 \\
循环动机 & riff & 短小、反复、作为记忆点的乐器句。 \\
钩子 & hook & 最容易记住的短句、音色或节奏。 \\
主调织体 & homophony & 一条主旋律加和弦伴奏，主次清楚。 \\
复调织体 & polyphony & 两条以上旋律同时独立进行。 \\
织体分层 & stratified texture & 低频、主体、空气、节奏各占一层。 \\
声部密度 & voice density & 同一时刻有多少条独立声部在说不同的话。 \\
音域 & range & 乐器或人声能发出的最低到最高音。 \\
移调 & transposition & 记谱音高与实际音高的换算，或整段旋律的挪动。 \\
适奏音区 & tessitura & 乐器/人声最舒服、最常待的一段音区。 \\
留白 & space & 故意不写满的空隙，让密度与对比可听。 \\
固定音型 & ostinato & 同一短句反复，变化靠位移、加花或换音色。 \\
升压段 & build-up & 电子舞曲进入 Drop 前的抬升段。 \\
电子落点 & drop & 电子舞曲中低频与主体同时回归的段落。 \\
侧链 & sidechain & 底鼓触发压缩，让 Pad/Bass 短暂让出低频。 \\
摇摆 & swing & 八分音符后半拍推迟形成的长—短律动。 \\
行走贝斯 & walking bass & 每拍一音、用经过音连接根音的贝斯线。 \\
强和弦 & power chord & 根音加五度的和弦，摇滚吉他常用。 \\
蓝调音 & blue note & 降三/降七并带弯曲的蓝调色彩音。 \\
影视配乐 & film score & 按画面与叙事组织段落、动机和配器的器乐。 \\
国风 & guofeng & 以五声音阶、民族音色、装饰与气口为特征的风格。 \\
情绪曲线 & emotional curve & 以段落为横轴，标出情绪强弱与释放点的走势。 \\
耳糖 & ear candy & 不承担主叙事、但让段落更好听的小音色或小句子。 \\
流行 & pop & 以钩子、律动、对比与兑现为核心传播逻辑的歌曲类型。 \\
托普线 & topline & 歌曲最上方的主旋律，通常由主唱承担。 \\
\bottomrule
\end{longtable}

\section*{本册不承诺的事}
\addcontentsline{toc}{section}{本册不承诺的事}

本册不替代混音、母带与插件说明书；不保证每种乐器组合都适合你的曲子；也不把“专业编曲”写成一套唯一正确的配方。所有音区、频率与力度数字都是\textbf{经验起点}，不是法律。先用它排除明显错误，再用耳朵和参考曲做最后判断。

\begin{helpmapping}
\begin{longtable}{p{0.45\textwidth}p{0.45\textwidth}}
\toprule
本节 & 帮助锚点 \\
\midrule
本册结构与学习路线 & \texttt{arrangement-plan} \\
十一章地图 & \texttt{arrangement-chapter-map} \\
三条学习路线 & \texttt{arrangement-learning-paths} \\
写作与验收约定 & \texttt{arrangement-conventions} \\
本册术语表 & \texttt{arrangement-glossary} \\
\bottomrule
\end{longtable}
\end{helpmapping}
